\section{PassageMeta.sty}\label{sec:ref-passagemeta}

Greppable \emph{passage} metadata --- the third and youngest of the Meta
packages, and the one that breaks the other two's rules. A passage is prose that
belongs to the author rather than to a course (see
section~\ref{sec:prose-documents}). Unlike its siblings this block is not a
search index: it is what gives \cs{PassageTitle} its heading, and what lets
\file{find-passages} write an import line for you. One command:

\begin{manexsource}
\PassageMeta{
  title    = {How We Mark},
  label    = {sec:howwemark},
  requires = {tikz},
}
\PassageTitle
\end{manexsource}

Keys: \env{title}, \env{label}, \env{requires}. Three, where a problem has six.
A passage has \emph{no} \env{source} or \env{number} (it is not drawn from a
textbook and it is not numbered), nothing extrinsic like points, and --- the part
worth explaining --- \textbf{no search keys at all}: see below.

\subsection*{The two rules it breaks}

\textbf{It stores.} ProblemMeta and FactMeta are store-only: every handler is
empty, the values discarded, the block present so that a mistyped key errors at
compile time. Here \env{title} and \env{label} have real handlers, because a
passage differs from a fact in \emph{where its name is consumed}. A fact's
\env{name} is consumed at the import site --- \cs{importfact}\texttt{[name=...]}
titles the box --- so \file{find-facts} injects it from the block into the call
site and FactMeta never needs the value itself. A passage's title is consumed
\emph{inside the passage file}, by \cs{PassageTitle}, a few lines below the
block. There is no import site to inject it into. So either the package stores
it, or the same string is written twice in one file --- and one job under two
spellings is what cost \env{desc} and \env{name} a rename
(section~\ref{sec:ref-problemmeta}).

\textbf{It is not removable.} Delete FactMeta and every fact still builds.
Delete this one and every passage loses its heading, so the failure is made loud
rather than quiet: \cs{PassageTitle} with no title raises a class error naming
what is missing. This is the one metadata package the toolkit cannot do without.

\subsection*{Nothing is injectable}

\texttt{find-passages} accepts no \texttt{-{}-with} at all, and it is the only
kind of which that is true. The one option an \cs{importpassage} line takes is
\env{level}, and \env{level} says where the passage's natural depth \emph{lands
in the document importing it} --- a property of that document, not of the
passage. Emitted lines are therefore always bare.

A \env{depth} key was tried here and removed, which is worth recording so it is
not re-derived. The idea was to note how far a passage reaches below its own
title, so you could tell before importing whether
\env{level}\texttt{=subsection} was legal. It gates nothing: the shift maps
\cs{section} to \cs{subsection} and \cs{subsection} to \cs{subsubsection}, the
class styles three ranks, so a title plus one rank of subsections still fits ---
importing the deepest passage in the bank at \env{level}\texttt{=subsection}
gives clean \texttt{1.1.1}'s. The limit bites only a passage carrying
\cs{subsubsection} itself, which none does, and \file{CourseDocument} already
refuses that at build time with a message naming the fix
(section~\ref{sec:ref-coursedocument}). Metadata duplicating a check the
compiler performs is a second reader of one fact.

\subsection*{Nothing to search by, deliberately}

The problem and fact banks are searched by \env{topics}; this one is searched by
nothing, and \file{find-passages} offers no filters. Two keys were tried and
withdrawn the day they arrived.

\env{topics} went first. A problem lives in \file{bacteria-doubling.tex}, whose
name says nothing about what it tests, so \env{topics} is the only way to ask
what is in the bank. A passage lives in \file{why-learn-math.tex}, whose name
\emph{is} its title --- and the title is a key here, which the picker
fuzzy-matches. Tagging the bank is what settled it: \env{marking} covered eight
of the eleven files while eight other values covered one file each.

\env{audience} (\env{students} / \env{tas} / \env{instructors}) went second,
and took the last filter with it. It had no build-time effect; its filter
duplicated the picker, which matches the displayed columns already; and it
restated the import graph, every passage being imported by exactly one document.
Its multi-valued tags were worse than redundant --- two passages marked for
students appear only in the markers' manual, so the tag recorded that
\emph{another} document holds similar prose rather than anything about the file.

What is left is the honest shape for a bank of dozens rather than hundreds, whose
files are named after their titles. Finding a passage was never the difficulty;
typing \file{../teaching/passages/} was, and \texttt{-{}-import -{}-from} is the
answer to that. Re-add \env{audience} if a passage is ever imported by two
documents --- the situation it was invented for, and the one that does not exist.

\subsection*{requires}

The one key with no analogue in the other two banks, and the reason a passage can
fail in a way a problem cannot: it lists what the \emph{importing document} must
supply. Course macros and packages both, because from the author's side they are
one question --- what must be in place before this passage compiles? --- even
though the fixes differ. The leading backslash tells them apart:

\begin{manexsource}
requires = {\MarkerCount}            % a macro the course must define
requires = {tikz}                    % a package the document must load
requires = {\MarkerCount, \PointsPerQuestion}
\end{manexsource}

\noindent Absent means self-contained; write nothing rather than
\env{requires}\texttt{ = \{\}}. There is no machinery behind it on purpose: a
course that imports a passage without supplying what it needs gets an
undefined-control-sequence error, which is the loud failure wanted. The key is in
\texttt{-{}-vocab} for the reason FactMeta's \env{kind} is, rather than the
reason a problem's \env{topics} is: it is functional, so \texttt{MarkerCount}
where the bank says \cs{MarkerCount} is a passage that will not build, not merely
a split search.

\subsection*{label, and what \LaTeX{} keeps}

\env{label} carries the \texttt{sec:} prefix every other sectioning label in the
courses uses, and is optional: omit it for a passage nothing refers to. It names
the passage's own idea, so it belongs to the passage rather than to whichever
document imports it.

On the \LaTeX{} side \env{requires} alone is \emph{validated but discarded},
exactly as in the siblings; \env{title} and \env{label} survive the block. The
real reader of all three is \file{find-passages}, which parses the
block textually off disk --- hence the same two disciplines (one key per line
with braced values, and \texttt{\#} written \texttt{\char`\\\#}) as the other two.
A macro name in \env{requires} needs no such care: the handler never expands it.
